\chapter{Introduction}
\label{chap:intro}

\section{Background and Context}

Losses during harvest are still one of the severe issues in the international food systems in
developing nations where the harvests efficiency has been significantly limited by the defects of
the supplying channel that have reportedly resulted in the shortage of food stocks and crop
harvests. With every harvest, a high percentage of crop products is not utilized; the level of
wastage is higher in places like Sub-Saharan Africa and South Asia due to lack of infrastructure
and lack of technology. These losses also weaken the economic sustenance of smallholder farm
structures besides depleting food security.

Farmer income is especially affected by the effects of post-harvest losses. Smallholder farmers,
who form a significant portion of agricultural producers in developing economies, tend to be
financially vulnerable. Losses realized after harvest directly reduce marketable output and
decrease profits. Research has indicated that handicapping in after-harvest processing and storage
considerably reduces economic returns, contributing to poverty cycles amongst agricultural
populations~\cite{opara2025application}. Uncertain losses also cause a rise in uncertainty,
making it difficult for farmers to make sound production and investment decisions.

Another big issue associated with post-harvest losses is food security. With world hunger on an
increasing curve, curtailing loss is an option to enhance food supply without necessarily
increasing food production. \cite{fadiji2023artificial} note that sustainable realization of food
systems requires the reduction of post-harvest losses in areas with already existing food
shortages. It is particularly troubling that resources spent on perishable goods that are not
consumed result in a loss of access to nutritious food as well as dietary deficits.

In many developing countries, supply chains tend to be inefficient and often fragmented, which
increases post-harvest losses. Inadequate transport systems, inefficient storage, and lack of
coordination between parties add to delays and rotting of products. In wholesale markets, fresh
produce is wasted in high quantities as a result of poor demand prediction and management
methods~\cite{opara2025application}. This underlines the necessity for better coordination and
the use of data-driven decision making across the supply chain.

The latest developments in artificial intelligence and machine learning present an opportunity to
tackle such issues. Predictive models have the ability to identify trends in environmental and
operating data that can predict any losses and help in making proactive decisions. The
implementation of neural network-based strategies to forecast post-harvest losses at production
and market scales has shown the potential of data-driven tools in enhancing
efficiency~\cite{guo2025predicting}. Likewise, predictive modeling in the garri production
demonstrated that machine learning can detect the main sources of losses and assist in targeted
changes~\cite{adewumi2025artificial}.

Although these advances have been made, the uptake of these technologies is still low in
smallholder situations because of financial, technical, and data availability constraints. This
highlights why basic, easy-to-reach machine learning solutions suited to the realities of the
developing world are necessary, which underpins the current research.


\section{Problem Statement}

In developing nations, post-harvest losses are still a major problem, with very high percentages
of agricultural products being wasted during storage, handling, and distribution. Research shows
that losses of up to 30 per cent are possible, especially in low-resource environments where
management practices and infrastructure are subpar~\cite{affognon2015unpacking}. The net effect
of these continuous losses is a weakening of food security, decreased incomes of farmers, and a
curtailment of agricultural systems on the whole.

Despite the potential of technological changes through artificial intelligence and machine learning
to overcome post-harvest inefficiencies, current solutions are not accessible to smallholder
farmers. Several of these systems demand sophisticated digital infrastructure, high-quality data,
and technical experience that are not suitable in resource-constrained settings. Moreover, such a
transformation in rural agricultural settings is not supported financially and is hampered by a
lack of digital literacy~\cite{wolfert2017big}.

Another equally important constraint is that the majority of currently available tools are geared
towards large-scale commercial farming with little regard given to the realities of running a
smallholder agribusiness. Consequently, this leads to an evident gap between technological
solutions and end-user needs. Also, existing methods tend to be either monitoring or
control-oriented, which restricts their ability to be proactive in preventing loss.

Therefore, there is an urgent need to develop simple, cheap, and interpretable predictive
instruments based on easily accessible data to forecast the risk of post-harvest losses. It is
important to bridge this gap and ensure that smallholder agricultural systems are resilient and
profitable through data-driven decision-making.


\section{Research Aim, Objectives and Questions}

\subsection{Aim}

The aim of this study is to identify and validate machine learning models capable of predicting
post-harvest loss risk using low-cost data.

\subsection{Research Objectives}

\begin{itemize}
    \item To determine and assess the storage, transportation, handling, and environmental
          conditions that contribute most to post-harvest losses.
    \item To assess the predictive accuracy of simple machine learning models by using sparse
          data to predict the post-harvest loss risk for smallholder crops.
    \item To evaluate the ways machine learning forecasts can be used to aid realistic
          agribusiness management choices.
\end{itemize}

\subsection{Research Questions}

\begin{enumerate}
    \item What storage, transportation, handling, and environmental factors have the greatest
          effect on post-harvest losses?
    \item To what extent can simple machine learning models correctly indicate whether a
          post-harvest loss is going to occur based on only limited data?
    \item What can these projections do to facilitate business management of agriculture?
\end{enumerate}


\section{Significance of the Study}

The research is significant at the academic level, the practice level, and at the policy level.
Academically, it promotes the use of machine learning in agricultural decision support systems,
especially in the under-researched field of post-harvest loss prediction. The study addresses the
gap between theoretical machine learning methods and practical agricultural problems by
emphasizing simple and understandable models.

In practice, the study provides a cheap and readily available solution to smallholder
agribusinesses. It offers a tool to help farmers predict the presence of post-harvest risks and
make informed decisions about storage, transportation, and handling by using readily available
data and computationally efficient models. This assists in minimizing losses and maximizing
profitability.

In policymaking, the findings promote long-term initiatives regarding improvement in food security
in developing nations. The study is strategic to long-term agricultural development and the
efficient use of resources, since it offers the capacity to employ data to conduct interventions
and enhance efficiency across the supply chain.


\section{Scope and Limitations}

The paper focuses on the development of a prototype machine learning model to forecast the hazard
of post-harvest losses in a smallholder agribusiness in the developing world. It is not a complete
research of all farm produce but is narrowed down to a small number of farm crops so that it can
be manageable and allow for an in-depth study.

The research involves simple, structured tabular data and does not involve complicated methods
like deep learning because such methods require large amounts of data and high computational
power, which are not available in low-resource settings. As an alternative, the focus is directed
toward interpretable and inexpensive models that can be implemented in practice.

There are a few limitations that are acknowledged. The data will be relatively small, potentially
impacting the model's predictability and generalisability. Also, the research is limited to
certain geographical and operational settings, which can restrict the extrapolation of results to
other areas. Being a prototype, the system is not implemented in real-time situations.


\section{Dissertation Structure}

There are six major chapters in this dissertation. Chapter~\ref{chap:intro} presents the research
topic and problem, and outlines the research objectives and questions. Chapter~\ref{chap:found}
provides a review of literature in the field of post-harvest losses and the application of machine
learning in agriculture. Chapter~\ref{chap:related} covers machine learning models for prediction,
decision support systems, the research gap, and the conceptual framework. Chapter~\ref{chap:method}
describes the research methodology, including data collection, model choice, and data evaluation
procedures. Chapter~\ref{chap:eval} reports the results and discussion of the machine learning
models. Lastly, Chapter~\ref{chap:concl} sums up the research, outlines the contributions, and
offers future research proposals.
